\documentclass[12pt]{article}
\usepackage[utf8]{inputenc}
\usepackage[T1]{fontenc}
\usepackage{hyperref}
\title{Tunen Language Resources for K-12 Curriculum Development}
\author{Compiled by Claude Sonnet 5 (claude-sonnet-5) for Andr\'as Kornai}
\date{}
\begin{document}
\maketitle

\section{Introduction}
Tunen (also Nen, Banen; ISO 639-3 tvu) is a Mbam Bantu language of Cameroon, spoken mainly in the Centre and Littoral regions. Webonary cites a population of over 100,000 for the Tunen group (1997), and lists several dialects. Wikipedia notes the language has unusual SOV word order for Bantu.

\section{Primary Educational Resources}
No school textbooks or primers were found in search results. Webonary mentions literature and texts in Tunen, but the search did not show specific titles.

\section{Linguistic and Reference Materials}
\begin{itemize}
\item Tunen--English dictionary on Webonary (Virginia Boyd, SIL International, 2019; a French version is also mentioned): \url{https://www.webonary.work/nen/overview/introduction/?lang=en}
\item Webonary overview: \url{https://www.webonary.work/nen/overview/?lang=en}
\item E.J. Kerr, \emph{Tunen syntax and information structure} (LOT dissertation, 2024), Leiden repository: \url{https://scholarlypublications.universiteitleiden.nl/access/item%3A4054928/view}
\item Wikipedia, Nen language (Cameroon): \url{https://en.wikipedia.org/wiki/Nen_language_(Cameroon)}
\end{itemize}

\section{Religious and Community Resources}
\begin{itemize}
\item Scripture Earth listing for Tunen (a 2017 CABTAL recording on Bible.is and a text on YouVersion): \url{https://www.scriptureearth.org/00eng.php?iso=tvu}
\end{itemize}
The extent of the scripture (complete Bible or portions) could not be confirmed.

\section{Digital and Audio Resources}
The Scripture Earth page mentions audio scripture and Global Recordings material, but no separate Global Recordings URL was seen in the results.

\section{Curriculum Development Context}
The language is not known to be taught in schools. The dictionary and the dissertation are good foundations for an orthography and lexicon. Graded readers and textbooks would still need to be created.
\end{document}